\section{Approach}

\subsection{Simulation Setup}
The simulation creates a controlled setting in which the true treatment effects are known. Each user $i$ has a group assignment $G_i$ and covariates $Z_i \in \mathbb{R}^5$. Daily wear is modeled as a two-state Markov chain, where state $1$ means the device is worn and state $0$ means it is not worn. The baseline transition probabilities are
\begin{align*}
    p_i &= P_i^{11} = \sigma(f_1(G_i, Z_i)), \\
    q_i &= P_i^{00} = \sigma(f_2(G_i, Z_i)),
\end{align*}
where $p_i$ is the probability of remaining compliant, $q_i$ is the probability of remaining noncompliant, and $\sigma$ is the sigmoid function.

The intervention effect also depends on $G_i$ and $Z_i$:
\begin{align*}
    \tau^p_i &= \sigma(g_1(G_i, Z_i)), \\
    \tau^q_i &= \sigma(g_2(G_i, Z_i)).
\end{align*}
Under treatment, the transition probabilities become
\begin{align*}
    p'_i &= p_i + (1-p_i)\tau^p_i, \\
    q'_i &= q_i(1-\tau^q_i).
\end{align*}
This parameterization keeps probabilities in $[0,1]$: treatment increases persistence in the compliant state and decreases persistence in the noncompliant state.

For each user, the initial state is sampled from the stationary distribution of the baseline Markov chain. We first simulate a 30-day pre-intervention history $H_i$ using $p_i$ and $q_i$. We then randomize treatment $A_i$ and simulate the next 30 days. Control users follow the baseline transition probabilities, while treated users follow $p'_i$ and $q'_i$. The post-intervention outcome $Y_i \in [0,1]$ is the fraction of compliant days in this second 30-day period. The simulated data are split into training and test sets.

\subsection{Two-Stage Predict-Then-Optimize}
The predict-then-optimize (PTO) pipeline first trains a predictive model and then solves a downstream optimization problem using the model predictions. The predictive model is trained on the training set with features $(X_i,A_i)$, where $X_i=(G_i,Z_i,H_i)$, and target $Y_i$. At decision time, the model is evaluated twice for each test user, once with $A_i=0$ and once with $A_i=1$. This produces estimated potential outcomes
\[
\hat{Y}_i^{(0)}, \qquad \hat{Y}_i^{(1)},
\]
which are used to estimate each user's predicted treatment benefit.

The estimated potential outcomes are then plugged into the objective in Equation~\ref{eq:obj}. The ranking setting corresponds to a nudging intervention, where each selected user consumes one unit of outreach capacity. The knapsack setting corresponds to a payment incentive, where each selected user has an estimated expected payment cost.

The PTO objective is
\begin{multline} \label{eq:obj_pto}
\max_A
\quad
\lambda_1 \sum_i A_i \left(\hat{Y}_i^{(1)} - \hat{Y}_i^{(0)}\right)
+ \lambda_2 \sum_i A_i \left(u(\hat{Y}_i^{(1)}) - u(\hat{Y}_i^{(0)})\right) \\
- \lambda_3 \sum_{g \in \mathcal{G}} \left(\hat{r}_g - \hat{r}_{\text{overall}}\right)^2 .
\end{multline}
For the ranking/nudging setting, the budget constraint is
\[
\sum_i A_i \le B.
\]
Without the fairness penalty, this reduces to sorting users by predicted objective contribution and selecting the top $B$ users. With the fairness penalty, the problem becomes a mixed-integer quadratic program (MIQP), which can be solved directly with Gurobi for the problem sizes considered here.

For the knapsack/incentive setting, the budget constraint is
\[
\sum_i \hat{C}_i A_i \le B,
\]
where the estimated payment cost is
\[
\hat{C}_i
=
\alpha \hat{Y}_i^{(1)}\mathbf{1}\{\hat{Y}_i^{(1)} \ge T\}.
\]
This produces a budgeted selection problem in which users are chosen based on both predicted benefit and predicted cost.

\subsection{Common DFL Score, Fairness, and Loss Structure}
Decision-focused learning (DFL) trains the predictive model through the downstream decision objective. The main difficulty is that the downstream decision rule contains non-differentiable operations, including thresholding, top-$B$ ranking, and greedy knapsack allocation. The DFL methods below differ in how they handle the non-differentiable allocation layer, but they share the same predicted outcomes, threshold scores, fairness computation, and full training loss.

For each user $i$, the model predicts the potential outcomes under no action and action:
\[
\hat{Y}_i^{(0)} = \sigma(f_\theta(X_i,0)),
\qquad
\hat{Y}_i^{(1)} = \sigma(f_\theta(X_i,1)).
\]
The hard threshold utility is approximated by a smooth sigmoid score
\[
\hat{S}_i^{(a)}
=
\sigma\!\left(
\frac{\hat{Y}_i^{(a)} - c}{\tau_{\text{thr}}}
\right),
\qquad a \in \{0,1\},
\]
where $c$ is the usable-data threshold and $\tau_{\text{thr}}>0$ controls the sharpness of the approximation. The marginal threshold score is
\[
d_i = \hat{S}_i^{(1)} - \hat{S}_i^{(0)}.
\]
In the ranking setting, $d_i$ is the score used to prioritize users. In the knapsack setting, the expected payment cost is approximated by
\[
\hat{c}_i
=
\alpha \hat{Y}_i^{(1)}
\sigma\!\left(
\frac{\hat{Y}_i^{(1)} - c_{\text{pay}}}{\tau_{\text{thr}}}
\right),
\]

All DFL methods are trained using stochastic mini-batches rather than the full population. In the ranking/nudging setting, the allocation constraint is naturally scaled to the mini-batch by selecting a fixed fraction of users. In the knapsack/incentive setting, however, the global budget constraint must be approximated at the batch level. We therefore define a mini-batch budget
\[
B_{\text{batch}}
=
\rho \sum_i \hat c_i,
\]
where the sum is taken over the current mini-batch. The parameter $\rho \in (0,1)$ controls the budget ratio. During training, $\rho$ may either be fixed or randomly sampled from a prespecified interval. Randomizing $\rho$ exposes the model to multiple budget regimes during training, while a fixed $\rho$ specializes the model to a single operating budget.

All DFL methods produce an allocation value $\tilde A_i \in [0,1]$ or a hard allocation value $A_i \in \{0,1\}$. Given an allocation vector $\bar A$, where $\bar A$ may denote a soft allocation, an averaged randomized allocation, or a deterministic hard allocation, the policy-level smoothed success value is
\[
\hat{S}_i^\pi
=
\bar A_i \hat{S}_i^{(1)}
+
(1-\bar A_i)\hat{S}_i^{(0)}.
\]
The shared fairness penalty is then
\[
\mathcal{L}_{\text{fair}}
=
\sum_{g\in\mathcal{G}}
\left(
\frac{1}{|G_g|}
\sum_{i\in G_g}
\hat{S}_i^\pi
-
\frac{1}{n}
\sum_i
\hat{S}_i^\pi
\right)^2.
\]
The full training loss combines a decision-focused loss, a supervised prediction loss, and the fairness penalty:
\[
\mathcal{L}
=
\lambda_{\text{DFL}}\mathcal{L}_{\text{DFL}}
+
\lambda_{\text{fair}}\mathcal{L}_{\text{fair}}
+
\lambda_{\text{mse}}\mathcal{L}_{\text{mse}}.
\]
The sections below only describe how each DFL method constructs the allocation layer and the corresponding decision-focused loss.

\subsection{End-to-End Differentiable Approximation for DFL}
The end-to-end differentiable approach replaces the non-differentiable allocation rule with a smooth surrogate so that gradients can flow through the full decision pipeline.

\paragraph{Ranking allocation.}
In the ranking/nudging setting, the hard top-$B$ rule selects the $B$ users with the largest marginal scores $d_i$. We approximate this rule with a soft allocation:
\[
\tilde{A}_i
=
\operatorname{clip}_{[0,1]}
\left(
B
\frac{\exp(d_i/\tau_{\text{alloc}})}
{\sum_j \exp(d_j/\tau_{\text{alloc}})}
\right),
\]
where $\tau_{\text{alloc}}>0$ controls the smoothness of the allocation. The decision-focused loss is
\[
\mathcal{L}_{\text{DFL-rank}}
=
-
\frac{1}{n}
\sum_i \tilde{A}_i d_i.
\]
For the shared fairness computation, we set $\bar A_i=\tilde A_i$.

\paragraph{Knapsack allocation.}
In the knapsack/incentive setting, the hard allocation problem is
\[
\max_{A_i\in\{0,1\}}
\sum_i A_i d_i
\quad
\text{s.t.}
\quad
\sum_i A_i \hat{c}_i \le B.
\]
We approximate the value-per-cost ranking by
\[
r_i = \frac{d_i}{\hat{c}_i+\epsilon},
\]
and assign soft weights
\[
w_i =
\frac{\exp(r_i/\tau_{\text{alloc}})}
{\sum_j \exp(r_j/\tau_{\text{alloc}})}.
\]
The allocation is scaled to approximately satisfy the budget:
\[
\tilde{A}_i
=
\operatorname{clip}_{[0,1]}
\left(
\frac{w_i}{\sum_j w_j \hat{c}_j}
B
\right).
\]
The decision-focused loss is
\[
\mathcal{L}_{\text{DFL-knapsack}}
=
-
\frac{1}{n}
\sum_i \tilde{A}_i d_i.
\]
Again, the shared fairness computation uses $\bar A_i=\tilde A_i$.

\subsection{DFL with Randomized Smoothing}
Randomized smoothing keeps the hard allocation rule but applies it repeatedly to randomly perturbed scores. The averaged selection frequency is then used as a differentiable surrogate in the loss. This preserves the structure of the downstream decision rule more closely than the fully soft allocation, while still producing stable training signals.

\paragraph{Ranking allocation.}
For the ranking/nudging setting, the hard allocation is
\[
A_i(d)
=
\mathbf{1}\{i\in\operatorname{TopB}(d)\}.
\]
To smooth this rule, we draw $M$ Gaussian perturbations
\[
\tilde d_i^{(m)} = d_i + \varepsilon_i^{(m)},
\qquad
\varepsilon_i^{(m)}\sim\mathcal{N}(0,\sigma^2),
\qquad
m=1,\dots,M.
\]
For each perturbed score vector, we recompute the hard top-$B$ allocation:
\[
A_i^{(m)}
=
\mathbf{1}\{i\in\operatorname{TopB}(\tilde d^{(m)})\}.
\]
The smoothed allocation is the average selection frequency:
\[
\bar A_i
=
\frac{1}{M}
\sum_{m=1}^M A_i^{(m)}.
\]
The decision-focused loss is
\[
\mathcal{L}_{\text{RS-rank}}
=
-
\frac{1}{n}
\sum_i \bar A_i d_i.
\]

 \paragraph{Knapsack allocation.}
For the knapsack/incentive setting, we use the shared mini-batch budget construction described above. For each score vector, the hard allocation is computed by a greedy knapsack rule: users are sorted by

Randomized smoothing applies the same greedy rule to perturbed scores. For each perturbation $m$, compute
\[
\tilde d_i^{(m)} = d_i + \varepsilon_i^{(m)},
\qquad
\varepsilon_i^{(m)}\sim\mathcal{N}(0,\sigma^2),
\]
and then
\[
A^{(m)}
=
\operatorname{GreedyKnapsack}(\tilde d^{(m)},\hat c,B_{\text{batch}}).
\]
The averaged allocation is
\[
\bar A_i
=
\frac{1}{M}
\sum_{m=1}^M A_i^{(m)},
\]
and the decision-focused loss is
\[
\mathcal{L}_{\text{RS-knapsack}}
=
-
\frac{1}{n}
\sum_i \bar A_i d_i.
\]

In both ranking and knapsack settings, gradients do not pass through the hard allocation solver itself. The averaged allocation $\bar A$ is treated as fixed within the mini-batch, and the model is updated through the score term $d_i$, the prediction loss, and the shared fairness penalty.

\subsection{DFL with Perturbation Gradient}
The perturbation-gradient method also preserves the hard downstream allocation rule. Instead of replacing the allocation with an averaged selection frequency, it estimates a surrogate gradient by measuring how the hard-allocation objective changes when the score vector is perturbed.

Let $A(d)$ denote the hard allocation produced from scores $d$. In the ranking/nudging setting, $A(d)$ is the top-$B$ allocation. In the knapsack/incentive setting, $A(d)$ is the greedy knapsack allocation using scores $d$, costs $\hat c$, and budget $B$. The hard decision objective is
\[
J(d)
=
\sum_i A_i(d)d_i.
\]
For perturbation sample $m$, draw
\[
\varepsilon^{(m)}\sim\mathcal{N}(0,I),
\]
and form
\[
d^{(m,+)} = d + h\varepsilon^{(m)},
\qquad
 d^{(m,-)} = d - h\varepsilon^{(m)},
\]
where $h$ is the perturbation scale. The hard allocation rule is recomputed under each perturbed score vector:
\[
A^{(m,+)} = A(d^{(m,+)}),
\qquad
A^{(m,-)} = A(d^{(m,-)}).
\]
In the implementation, the perturbed allocations are evaluated using the original detached score vector:
\[
J^{(m,+)} = \sum_i A_i^{(m,+)}d_i,
\qquad
J^{(m,-)} = \sum_i A_i^{(m,-)}d_i,
\]
with baseline objective
\[
J(d)=\sum_i A_i(d)d_i.
\]

The perturbation-gradient estimator can be computed in three ways.

\paragraph{Forward estimator.}
The forward estimator compares the positively perturbed allocation with the baseline allocation:
\[
\widehat{\nabla}_d J(d)
=
\frac{1}{M}
\sum_{m=1}^M
\frac{J^{(m,+)}-J(d)}{h}
\varepsilon^{(m)}.
\]

\paragraph{Backward estimator.}
The backward estimator compares the baseline allocation with the negatively perturbed allocation:
\[
\widehat{\nabla}_d J(d)
=
\frac{1}{M}
\sum_{m=1}^M
\frac{J(d)-J^{(m,-)}}{h}
\varepsilon^{(m)}.
\]

\paragraph{Central estimator.}
The central estimator uses a symmetric finite-difference comparison:
\[
\widehat{\nabla}_d J(d)
=
\frac{1}{M}
\sum_{m=1}^M
\frac{J^{(m,+)}-J^{(m,-)}}{2h}
\varepsilon^{(m)}.
\]

The estimated gradient is treated as fixed during backpropagation and is used to define the surrogate decision-focused loss
\[
\mathcal{L}_{\text{PG}}
=
-
\frac{1}{n}
\sum_i
\widehat{\nabla}_{d_i}J(d)\, d_i.
\]
For the shared fairness computation, perturbation-gradient DFL uses the deterministic hard allocation from the current scores, $\bar A_i=A_i(d)$. Thus, the method does not differentiate directly through the top-$B$ or greedy knapsack solver; instead, it updates the score model in the direction suggested by randomized finite-difference perturbations of the hard decision objective.